% =====================================================================
% fig-2-8-workflow-skill-cache.tex  (version simplifiée)
%
% Le workflow de la phase 3 = celui de la phase 2 (fig-2-3-workflow-mcp),
% repris tel quel : harnais du chapitre 1 + client MCP + boucle d'outils
% (huit tours au plus). SEUL AJOUT dessiné : les deux points de coupe du
% cache de prompt (en ocre). Le skill n'apparaît pas : il ne change pas
% la forme du workflow, son corps entre par le prompt système dans
% « Préparer l'appel ».
%
% Bibliothèques TikZ requises : \usetikzlibrary{arrows.meta}
% Largeur du tracé : environ 15,5 cm.
% =====================================================================

\definecolor{wfbord}{HTML}{4A5568}
\definecolor{wffond}{HTML}{F7F7F4}
\definecolor{wfmcpb}{HTML}{2F6F66}
\definecolor{wfmcpf}{HTML}{E8F1EF}
\definecolor{wfcacheb}{HTML}{B3541E}
\definecolor{wfcachef}{HTML}{FBEDE4}
\definecolor{wffleche}{HTML}{6B7280}

\begin{tikzpicture}[
  font=\scriptsize,
  >={Stealth[length=2mm]},
  bloc/.style  = {draw=wfbord, fill=wffond, rounded corners=2pt,
                  line width=.4pt, text width=2.35cm, align=center,
                  inner sep=3pt, minimum height=10mm},
  mcp/.style   = {bloc, draw=wfmcpb, fill=wfmcpf},
  test/.style  = {bloc, dash pattern=on 2pt off 1.5pt},
  cache/.style = {draw=wfcacheb, fill=wfcachef, rounded corners=2pt,
                  line width=.8pt, text width=3.1cm, align=center,
                  inner sep=3pt, font=\tiny, text=wfcacheb},
  fl/.style    = {->, draw=wffleche, line width=.5pt},
  flc/.style   = {->, draw=wfcacheb, line width=.7pt,
                  dash pattern=on 2pt off 1.5pt},
  et/.style    = {font=\tiny, text=wffleche, inner sep=1.5pt, fill=white}
]

% ------------------------------------------------ chaîne principale
\node[bloc] (bcl) at ( 1.40, 0)    {Boucle\\[-2pt]{\tiny 42 $\times$ 3 = 126, un appel à la fois}};
\node[bloc] (prp) at ( 4.75, 0)    {Préparer l'appel\\[-2pt]{\tiny consigne + outils + journal}};
\node[bloc] (llm) at ( 8.10, 0)    {Appel LLM\\[-2pt]{\tiny Claude Opus 4.8}};
\node[test] (tmc) at (11.45, 0)    {Tour MCP ?\\[-2pt]{\tiny 8 tours au plus}};
\node[bloc] (jgm) at (14.80, 0)    {Jugement,\\[-3pt]ligne JSONL};

\draw[fl] (bcl) -- (prp);
\draw[fl] (prp) -- (llm);
\draw[fl] (llm) -- (tmc);
\draw[fl] (tmc) -- node[et, above] {triplet} (jgm);

% ------------------------------------------------ boucle d'outils
\node[bloc] (lot) at (11.45,-2.55) {Préparer le lot\\[-2pt]{\tiny un envoi JSON-RPC}};
\node[mcp]  (srv) at ( 8.10,-2.55) {Serveur MCP\\[-2pt]{\tiny \texttt{POST /mcp}}};
\node[mcp]  (itg) at ( 4.75,-2.55) {Intégrer\\[-2pt]{\tiny \texttt{tool\_result}}};

\draw[fl] (tmc.south) -- node[et, right] {outils demandés} (lot.north);
\draw[fl] (lot) -- (srv);
\draw[fl] (srv) -- (itg);
\draw[fl] (itg.north) -- ++(0,0.85) -| node[et, pos=.25, above] {l'appel repart} (llm.south);
\draw[fl] (jgm.south) -- (14.80,-3.85) -- (1.40,-3.85)
          node[et, above, pos=.55] {exécution suivante}
          -- (bcl.south);

% ------------------------------------------------ le seul ajout : le cache
\node[cache] (c1) at (4.75, 1.75) {\textbf{AJOUT PHASE 3 --- CACHE}\\point de coupe explicite en fin de bloc système, couvrant les outils};
\node[cache] (c2) at (1.60,-2.55) {\textbf{AJOUT PHASE 3 --- CACHE}\\cache automatique de tête pour la conversation qui grandit};
\draw[flc] (c1.south) -- (prp.north);
\draw[flc] (c2.east) -- (itg.west);

\end{tikzpicture}
